Explain in the Gaussian case how to obtain the distribution of Student’ized resid-
uals can be derived? What is the problem with a simple standardization?

Under the assumptions on the residuals, namely

1. Expectation: \(\mathbb{E}(\hat{\epsilon}) = 0\)
2. Variance: \(\text{Var}(\hat{\epsilon}_i) = \sigma^2 (1 - h_{ii})\) where \(h_{ii}\) is the \(i\)-th diagonal entry of \(H = \mathbf{X}(\mathbf{X}^\top \mathbf{X})^{-1} \mathbf{X}^\top\).
3. Covariance: \(\text{Cov}(\hat{\epsilon}) = \sigma^2 (I - H)\)

With normally distributed errors

1. Distribution of Residuals: \(\hat{\epsilon} \sim N(0, \sigma^2 (I - H))\)
2. Distribution of Residual Sum of Squares: \(\frac{\hat{\epsilon}^\top \hat{\epsilon}}{\sigma^2} = (n - p) \frac{\hat{\sigma}^2}{\sigma^2} \sim \chi^2_{n - p}\)
3. Independence: \(\hat{\epsilon}^\top \hat{\epsilon} = (n - p) \hat{\sigma}^2\) and \(\hat{\beta}\) are independent.

In order to remove heteroscedasticity, residuals may be standardized (simple approach): 

\( r_i = \frac{\hat{\epsilon}_i}{\hat{\sigma} \sqrt{1 - h_{ii}}} \).

However, since numerator and denominator are dependent, an analysis of the distribution is not easy.

An alternative are Student'ized residuals with a different estimator for \(\sigma\):

Suppose that errors are all normally distributed.
Removing the \(i\)-th response variable and covariate data, we obtain \(\mathbf{y}_{(i)}\) and \(\mathbf{X}_{(i)}\).
The corresponding least squares estimator (LSE) and estimator for \(\sigma^2\) that neglect the \(i\)-th observation are

\[
\hat{\beta}_{(i)} = \left( \mathbf{X}_{(i)}^\top \mathbf{X}_{(i)} \right)^{-1} \mathbf{X}_{(i)}^\top \mathbf{y}_{(i)}
\]
\[
\hat{\sigma}^2_{(i)} = \frac{1}{n - p - 1} \sum_{j = 1, 2, \ldots, n, \, j \neq i} \left( y_j - \mathbf{x}_j^\top \hat{\beta}_{(i)} \right)^2.
\]

In particular,
\[
\hat{\sigma}^2_{(i)} \quad \perp \quad \hat{\epsilon}_{(i)} := y_i - \mathbf{x}_i^\top \hat{\beta}_{(i)}.
\]

Using independence and 
\(\hat{\epsilon}_{(i)} \sim N \left(0, \sigma^2 \left(1 + \mathbf{x}_i^\top \left( \mathbf{X}_{(i)}^\top \mathbf{X}_{(i)} \right)^{-1} \mathbf{x}_i \right)\right)\),

\(
\left( n - p - 1 \right) \frac{\hat{\sigma}^2_{(i)}}{\sigma^2} \sim \chi^2_{n - p - 1},
\)
we get the Student'ized residuals
\[
r_i^* = \frac{\hat{\epsilon}_{(i)}}{\hat{\sigma}_{(i)} \sqrt{1 + \mathbf{x}_i^\top \left( \mathbf{X}_{(i)}^\top \mathbf{X}_{(i)} \right)^{-1} \mathbf{x}_i}} \sim t_{n - p - 1}.
\]


In the context of hypothesis testing in the Gaussian linear model. 
What are general linear hypotheses? 

General linear hypotheses are specified as follows:
\[ H_0: C\beta = d \quad \text{vs} \quad H_1: C\beta \neq d \]

with \( \mathbf{C} \in \mathbb{R}^{r, p} \) with \( \text{rk}(\mathbf{C}) = r \leq p \), i.e., \( r \) linearly independent conditions are tested.

Test of significance: 
\( H_0: \beta_j = 0 \quad \text{vs} \quad H_1: \beta_j \neq 0 \),
with corresponding
\( \mathbf{C}^\top = (\delta_{i, j+1})_i \)

Composite test:
\( H_0: \beta_1 = 0 \quad \text{vs} \quad H_1: \beta_1 \neq 0, \, \beta_1 = (\beta_1, \beta_2, \ldots, \beta_r)^\top \)

\[ \mathbf{C} = \begin{pmatrix}
0 & 1 & 0 & \cdots & 0 & \cdots & 0 \\
0 & 0 & 1 & \cdots & 0 & \cdots & 0 \\
\vdots & \vdots & \ddots & \vdots & 0 \\
0 & 0 & 0 & \cdots & 1 & 0 \\
\end{pmatrix} \]

Test of equality:
\( H_0: \beta_j - \beta_r = 0 \quad \text{vs} \quad H_1: \beta_j - \beta_r \neq 0 \)

\[ \mathbf{C}^\top = (\delta_{i, j+1} - \delta_{i, r+1})_i \]

A general drawback is that the LSE error is larger.



In the context of hypothesis testing in the Gaussian linear model.
Describe statistics to construct tests, e.g. \( SSE \)
and the corresponding distribution.

The least squares estimator under the linear constraint \(C \beta = d\) is given by

\(
\hat{\beta}^R = \hat{\beta} - (\mathbf{X}^\top \mathbf{X})^{-1} \mathbf{C}^\top \left( \mathbf{C} (\mathbf{X}^\top \mathbf{X})^{-1} \mathbf{C}^\top \right)^{-1} ( \mathbf{C} \hat{\beta} - d )
\)

The unrestricted residual sum of squares is

\(
SSE = \hat{\epsilon}^\top \hat{\epsilon} = (n - p) \hat{\sigma}^2
\)

The residual sum of squares under the hypothesis \(H_0\) equals
\(
SSE + \Delta SSE
\)
with

\(
\Delta SSE = (\mathbf{C} \hat{\beta} - d)^\top \left( \mathbf{C} (\mathbf{X}^\top \mathbf{X})^{-1} \mathbf{C}^\top \right)^{-1} ( \mathbf{C} \hat{\beta} - d )
\)
the test statistic is then \( \frac{\delta \text{SSE}}{\text{SSE}} \)
and under the null hypothesis \( H_0 \), it is additionally known that

\(
\frac{1}{\sigma^2} \Delta SSE \sim \chi^2_r
\)

We have already seen that, in any case, \(\frac{1}{\sigma^2} SSE \sim \chi^2_{n - p}\). 
This implies, due to independence, that under the hypothesis \( H_0 \):

\(
F := \frac{n - p}{r} \cdot \frac{\Delta SSE}{SSE} \sim F_{r, n - p}
\)
called an \( F \)-distribution.


The hypothesis \( H_0 \) is rejected, if \( F \) is large:
Since the distribution of \( F \) under \( H_0 \) is known, we obtain a test of level \(\alpha \in (0, 1)\) by rejecting \( H_0 \), if
\(
F > F_{r, n - p}(1 - \alpha).
\)
Here, \( F_{r, n - p}(1 - \alpha) \) denotes the \((1 - \alpha)\)-quantile of an \( F \)-distribution with \( r \) and \( n - p \) degrees of freedom.


In the context of hypothesis testing in the Gaussian linear model.
How can one use \( \frac{\Delta \text{SSE}}{\text{SSE}} \) this to set up a test of a given size? 

It simplifies under some assumptions,
given the estimators 

\( \hat{\text{Cov}(\mathbb{\beta})} = \hat{\sigma}^2 (\mathbb{X}^T\mathbb{X})^{-1}\)
and \( \hat{\sigma}^2 = \frac{1}{n - p}\varepsilon^T\varepsilon \)

denote as the diagonal entries of the covariance \( \text{se}_j \coloneqq \sqrt{\hat{\text{Cov}(\mathbb{\beta})}}\)

Tests at level \(\alpha\) can be constructed then as follows:

1. Test of significance \(H_0 : \beta_j = 0\) vs \(H_1 : \beta_j \neq 0\)

Under \(H_0\), the test statistic \(t_j = \frac{\hat{\beta}_j}{se_j}\) is \(t_{n - p}\)-distributed, and the null hypothesis \(H_0\) is rejected, if \(|t_j| > t_{n - p}(1 - \alpha/2)\) where \(t_{n - p}(1 - \alpha/2)\) denotes the \((1 - \alpha/2)\)-quantile of the \(t\)-distribution with \(n - p\) degrees of freedom.

2. Composite test \(H_0 : \beta_1 = 0\) vs \(H_1 : \beta_1 \neq 0\)

Under \(H_0\), the test statistic \(F = \frac{1}{r} \hat{\beta}_1^\top \text{Cov}(\hat{\beta}_1)^{-1} \hat{\beta}_1\) is \(F_{r, n - p}\)-distributed, and the null hypothesis \(H_0\) is rejected, if \(F > F_{r, n - p}(1 - \alpha)\) where \(F_{r, n - p}(1 - \alpha)\) denotes the \((1 - \alpha)\)-quantile of the \(F\)-distribution with \(r\) and \(n - p\) degrees of freedom.

3. Test of equality \(H_0 : \beta_1 = \beta_2 = \cdots = \beta_k = 0\) vs \(H_1 : \text{at least one} \ \beta_1, \ldots, \beta_k \neq 0\)

Under \(H_0\), the test statistic \(F = \frac{n - p}{k} \cdot \frac{R^2}{1 - R^2}\) is \(F_{k, n - p}\)-distributed, and the null hypothesis \(H_0\) is rejected, if \(F > F_{k, n - p}(1 - \alpha)\) where \(F_{k, n - p}(1 - \alpha)\) denotes the \((1 - \alpha)\)-quantile of the \(F\)-distribution with \(k\) and \(n - p\) degrees of freedom.


Applicable distributions are the Gamma, \(\chi^2\), Beta, Fisher and Student's t-distribution


The following results hold for normally distributed errors and asymptotically for large sample sizes.

1. Confidence interval for \(\beta_j\)
A confidence interval for \(\beta_j\) of level \(1 - \alpha\) is given by

\(
[\hat{\beta}_j - t_{n - p}(1 - \alpha/2) \cdot se_j, \hat{\beta}_j + t_{n - p}(1 - \alpha/2) \cdot se_j]
\)

2. Confidence ellipsoid for subvector \(\beta_1\)
A confidence ellipsoid for \(\beta_1 = (\beta_1, \beta_2, \ldots, \beta_r)^\top\) of level \(1 - \alpha\) is given by

\(
\left\{ \beta_1 : \frac{1}{r} (\hat{\beta}_1 - \beta_1)^\top \widehat{\text{Cov}}(\hat{\beta}_1)^{-1} (\hat{\beta}_1 - \beta_1) \leq F_{r, n - p}(1 - \alpha) \right\}
\)

In the context of hypothesis testing in the Gaussian linear model.
What are the prediction/confidence intervals for the future observation \( \hat{y}_0 \)

Assuming Gaussian errors, a confidence interval for \(\mu_0 = \mathbb{E}(y_0)\) of a future observation \(y_0\) 
at location \(\mathbf{x}_0\) at level \(1 - \alpha\) is given by

\[
\mathbf{x}_0^\top \hat{\beta} \pm t_{n - p}(1 - \alpha/2) \hat{\sigma} \sqrt{\mathbf{x}_0^\top (\mathbf{X}^\top \mathbf{X})^{-1} \mathbf{x}_0}
\]

derived from \( \mathbb{x}^T_0\hat{\mathbb{\beta}} \sim \mathcal{N}(\mathbb{x}^T_0\hat{\mathbb{\beta}} = \mu_0, \sigma^2 \mathbb{x}_0^T(\mathbb{X}^T\mathbb{X})^{-1}\mathbb{x}_0 \)
where \( \hat{\sigma} \sim \chi^2\) is used for \( \sigma \).


Closely related is the issue of a prediction interval for \( y_0 \), i.e., an estimated interval such that (independently of the true value of the parameters \(\beta\)) the realization of \( y_0 \) falls into this interval with probability at least \( 1 - \alpha \).
This question can be solved as follows:

The prediction error 
\[ \hat{\epsilon}_0 = y_0 - \mathbf{x}_0^\top \hat{\beta} = \mathbf{x}_0^\top \beta + \epsilon_0 - \mathbf{x}_0^\top \hat{\beta} = \mu_0 + \epsilon_0 - \mathbf{x}_0^\top \hat{\beta}\] 
has mean 0, and the summands \( \epsilon_0 \) and \( \mathbf{x}_0^\top \hat{\beta} \) are normally distributed and independent. Thus,

\(
\hat{\epsilon}_0 \sim N \left( 0, \sigma^2 + \sigma^2 \mathbf{x}_0^\top (\mathbf{X}^\top \mathbf{X})^{-1} \mathbf{x}_0 \right)
\)

Since \( \hat{\sigma}^2 \) is both independent of \( \epsilon_0 \) and \(\hat{\beta}\), it is independent of \( \hat{\epsilon}_0 \), implying that

\(
T := \frac{\hat{\epsilon}_0}{\hat{\sigma} \sqrt{1 + \mathbf{x}_0^\top (\mathbf{X}^\top \mathbf{X})^{-1} \mathbf{x}_0}} \sim t_{n-p}
\)

Obviously, \( P(|T| \leq t_{n-p}(1 - \alpha/2)) = 1 - \alpha \). Solving for bounds on \( y_0 \) produces the sought prediction interval,
wich is
\[
\mathbf{x}_0^\top \hat{\beta} \pm t_{n-p} \cdot \hat{\sigma} \sqrt{1 + \mathbf{x}_0^\top (\mathbf{X}^\top \mathbf{X})^{-1} \mathbf{x}_0}
\]
for a future observation \( y_0 \) at location \( \mathbf{x}_0 \) at level \( 1 - \alpha \)
